% Having a shared abstract here because before the main and submission versions
% of the document were not in sync. This should avoid this issue in the future.
\begin{abstract}

Population size is a key factor underlying the mode and tempo of evolution,
particularly as it relates to the strength of selection and drift.
While the mechanisms underlying the interactions between population size
and selection have been studied in population genetics for over a century, empirical
knowledge of these dynamics has been limited to species with small-to-moderate
historical population sizes. This gap in knowledge highlights the need for empirical
studies in systems with historically large population sizes and elevated diversity.
The Pacific acorn barnacle (\balgla) presents an exceptional model to study evolution
in extremely large populations, exhibiting census sizes often exceeding the tens of
thousands of individuals per square meter.
We generated a new chromosome-level genome assembly for this species, which we leverage
in one of the first large-scale genomic analyses in this system.
This assembly reveals a highly polymorphic genome with over 3\% heterozygosity.
At a population level, \bgla{} exhibits extreme levels of polymorphism,
with nucleotide diversity surpassing 5\% genome-wide in just a small collection of
individuals.
Across the genome, nucleotide diversity predictably decreases at functional elements,
including both coding and non-coding sequences, likely reflecting strong purifying
selection along the genome.
At the same time, McDonald-Kreitman tests reveal that the majority of non-synonymous
substitutions between barnacle species were driven by positive selection, consistent
with the expected increase in the efficacy of selection in large populations.
These remarkable levels of diversity set \bgla{} as a unique model for the study of
evolution at extreme demographic scales and highlights the importance of testing
evolutionary theory across a wide variety of empirical systems.

\end{abstract}
